% ============================================================================
\section{Local Fitting and Efficient Allocation}
\label{sec:theory}
% ============================================================================

We now introduce the checkpoint-level analysis evaluated in
Section~\ref{sec:experiments}. The analysis has two parts. First, a local
sampled-token fit separates solo reachability, shared-update conflict, and
finite-update optimization error, then predicts the response to task weights.
Second, a fixed-objective sampling analysis produces the optimizer-aware
allocation used by GPAS.


\subsection{A local quadratic model of MOPD}
\label{sec:local_fit}

Fix a checkpoint $\theta_0$. Let $\nu_i^{\theta_0}$ denote the distribution
of valid response-token pairs $z=(c,y)$ obtained from task-$i$ student
rollouts at that checkpoint. For one such pair, define the sampled-token
teacher signal and its score gradient,
\begin{equation}
    d_i(z)
    =\log\tau_i(y\mid c)-\log\pi_{\theta_0}(y\mid c),
    \qquad
    j_i(z)=\nabla_\theta\log\pi_{\theta_0}(y\mid c).
    \label{eq:local_logprob_terms}
\end{equation}
For compactness, let $J_i\Delta$ denote the scalar function
$z\mapsto j_i(z)^\top\Delta$, and define
\begin{equation}
    \|h\|_i^2
    =\mathbb E_{z\sim\nu_i^{\theta_0}}[h(z)^2].
    \label{eq:local_token_norm}
\end{equation}
In experiments this expectation is a mean over cached tokens. $J_i$ is only
a matrix-free linear map evaluated with Jacobian--vector and
vector--Jacobian products; no token-by-parameter Jacobian is stored.

Linearizing the sampled student log probability on the fixed cache gives
$\log\pi_{\theta_0+\Delta}(y\mid c)
=\log\pi_{\theta_0}(y\mid c)+j_i(z)^\top\Delta
+O(\|\Delta\|^2)$. We therefore use the empirical Gauss--Newton surrogate
\begin{equation}
    L_w(\Delta)
    =
    \frac12\sum_{i=1}^{M}w_i
    \left\|d_i-J_i\Delta\right\|_i^2.
    \label{eq:local_projection}
\end{equation}
This model fits the same sampled-token log-ratio signal used by training. It
is a checkpoint-level empirical diagnostic, not an exact representation of
the full-vocabulary KL or a global model of the training trajectory.
Appendix~\ref{app:local_kl} gives its construction and its local connection
to reverse KL.


\subsection{Fitting-Gap Decomposition}
\label{sec:mixture_sensitivity}

Let $\Delta_w^*$ minimize the unconstrained local surrogate in
Equation~\ref{eq:local_projection}, and let $L_w^*=L_w(\Delta_w^*)$.
The weighted least-squares geometry gives the following decomposition.

\begin{proposition}[Local fitting and optimization error]
\label{prop:gap_decomposition}
For every local update $\Delta$,
\begin{equation}
    L_w(\Delta)
    =
    L_w^*
    +
    \frac12\sum_{i=1}^{M}w_i
    \left\|J_i(\Delta-\Delta_w^*)\right\|_i^2.
    \label{eq:gap_decomposition}
\end{equation}
The first term is the residual of the best shared update within this local
model. The second is the additional local error from not reaching that
update.
\end{proposition}

Appendix~\ref{app:projection} proves the proposition. The residual $L_w^*$
alone does not identify cross-task conflict: it can remain positive even for
one task whose requested change is outside the student's local tangent space.
To isolate the cost of sharing an update, define
\begin{align}
    \ell_i(\Delta)
    &:=\frac12\|d_i-J_i\Delta\|_i^2,\\
    L_{i,\mathrm{solo}}^*(R)
    &:=\min_{\|\Delta_i\|_2\le R}\ell_i(\Delta_i),
    \qquad
    L_{w,\mathrm{solo}}^*(R)
    :=\sum_iw_iL_{i,\mathrm{solo}}^*(R),\\
    L_w^*(R)
    &:=\min_{\|\Delta\|_2\le R}\sum_iw_i\ell_i(\Delta).
    \label{eq:solo_joint_fit}
\end{align}
The joint optimum obeys
$L_w^*(R)\ge L_{w,\mathrm{solo}}^*(R)$ because separate fits can choose a
different update for each task. For $L_w(0)>0$ and
$\|\Delta_{\rm train}\|_2\le R$, the gap diagnostic is
\begin{align}
    \mathrm{SoloGap}(R)
    &:=\frac{L_{w,\mathrm{solo}}^*(R)}{L_w(0)},\\
    \mathrm{ConflictGap}(R)
    &:=\frac{L_w^*(R)-L_{w,\mathrm{solo}}^*(R)}{L_w(0)},\\
    \mathrm{FitGap}(R)
    &:=\frac{L_w^*(R)}{L_w(0)}
    =\mathrm{SoloGap}(R)+\mathrm{ConflictGap}(R),\\
    \mathrm{OptGap}(R)
    &:=\frac{L_w(\Delta_{\rm train})-L_w^*(R)}{L_w(0)}.
    \label{eq:gap_diagnostics}
\end{align}
These quantities satisfy
$\mathrm{SoloGap}(R)+\mathrm{ConflictGap}(R)+\mathrm{OptGap}(R)
=L_w(\Delta_{\rm train})/L_w(0)$ and are nonnegative. The solo gap measures
residual fit error when each task gets its own update in the tested local
model and radius. It is not a claim about global student capacity. The
conflict gap is the additional local residual caused by requiring one shared
update. These quantities depend on the cache, parameterization, and radius;
we therefore report held-out estimates, solver diagnostics, and constraint
status. Exact nonnegativity holds for the objective on which the solutions
are defined; cross-fitted estimates can deviate slightly because of finite
cache error.


\subsection{Prediction for Task-Weight Interventions}
\label{sec:weight_response}

The residual in the best local fit also determines how a change in the target
mixture moves each task. Let $\bar\Delta_w$ denote the joint solution used for
the diagnostic: $\Delta_w^*$ when the solution is interior, and
$\arg\min_{\|\Delta\|_2\le R}L_w(\Delta)$ when the trust-region boundary is
active. Define
\begin{equation}
    H_w=\sum_i w_i
        \mathbb E_i[j_i(z)j_i(z)^\top],
    \qquad
    r_i=\mathbb E_i
        [j_i(z)(d_i(z)-j_i(z)^\top\bar\Delta_w)],
    \label{eq:local_curvature_simple}
\end{equation}
where $\mathbb E_i$ is over $z\sim\nu_i^{\theta_0}$. Here $H_w$ denotes a
positive-semidefinite linear operator, not a matrix that is formed or
inverted. Parameterize the weights as $w=\operatorname{softmax}(\eta)$ and let
$u_i=\partial\bar\Delta_w/\partial\eta_i$. For an interior solution, the
response is obtained from the matrix-free linear solve
\begin{equation}
    H_wu_i=w_i r_i.
    \label{eq:weight_sensitivity}
\end{equation}
For a damped interior solution, the implementation instead solves
$(H_w+\xi I)u_i=w_i(r_i-\xi\bar\Delta_w)$. If the radius is active, we
differentiate the trust-region optimality conditions. All systems are solved
with matrix-free conjugate gradients; Appendix~\ref{app:projection} gives the
details.
For the local task loss
$\ell_j(\Delta)=\tfrac12\|d_j-J_j\Delta\|_j^2$, a small intervention
$\delta\eta_i$ therefore gives
\begin{equation}
    \delta \ell_j
    \approx
    -r_j^\top u_i\,\delta\eta_i.
    \label{eq:task_loss_response}
\end{equation}
The cross-task inner products in Equation~\ref{eq:task_loss_response} predict
which other sampled-token teacher losses improve or deteriorate.
Section~\ref{sec:exp_mixture} uses the interior or active-constraint solve
according to the reported constraint status and tests the full response
matrix with matched short branches.


\subsection{Optimizer-Aware Allocation Rule}
\label{sec:gpas}

The preceding analysis concerns the target mixture $w$. We next hold that
mixture fixed and choose the sampling probabilities $a$. Before sampling at
step $t$, freeze AdamW's lagged second-moment preconditioner,
\begin{equation}
    P_t=\operatorname{diag}
    \left((\sqrt{v_{t-1}}+\epsilon)^{-1}\right),
    \qquad
    U_i=P_tG_i,
    \label{eq:optimizer_update}
\end{equation}
and define the task signal
\begin{equation}
    s_i=\sqrt{\mathbb E\|U_i\|_2^2}.
    \label{eq:optimizer_scale}
\end{equation}
This score is a lagged-preconditioned gradient norm, used as a proxy for the
task-dependent AdamW step scale. It is not the exact AdamW update. The
first-order loss decrease also depends on gradient direction, so we treat loss
reduction as a separate empirical outcome.

For one sampled task and fixed $P_t$, minimizing the second moment of the
importance-corrected estimator gives
\begin{equation}
    a_i^*
    =
    \frac{w_i s_i}{\sum_jw_js_j}.
    \label{eq:optimal_sampling}
\end{equation}
Because Equation~\ref{eq:importance_unbiased} fixes the conditional mean,
this also minimizes
$\mathbb E\|P_t\widehat g-P_tg_w\|_2^2$. The importance-sampling
minimization is classical; the modeling choice here is to apply it at the
fixed-token task level using the lagged AdamW proxy. The short proof is
in Appendix~\ref{app:optimal_sampling}.

Gradient-Preserving Adaptive Sampling (GPAS) estimates each score online:
\begin{equation}
    \widehat s_i^2
    \leftarrow
    \gamma\widehat s_i^2
    +(1-\gamma)\|P_tG_i\|_2^2
    \label{eq:gpas_ema}
\end{equation}
after sampling task $i$. A probability floor limits stale statistics and
large importance ratios. The exact lower-bounded solution is
\begin{equation}
    a_i
    =\max\left\{a_{\min},
    c\,w_i\sqrt{\widehat s_i^2+\varepsilon}\right\},
    \qquad \sum_i a_i=1,
    \label{eq:gpas_distribution}
\end{equation}
where $c$ enforces normalization. The sampled gradient is multiplied by
$w_I/a_I$ before the AdamW moment update. The name refers specifically to
preserving the conditional mean of the target gradient estimator. GPAS does
not preserve the expected AdamW parameter update or the optimizer trajectory.
It stores one score and one timestamp per task and requires no extra model
evaluation.

Raw GPAS sets $P_t=I$. Adam GPAS uses Equation~\ref{eq:optimizer_update}.
The latter relies on the quality of the lagged approximation, which we test
in Section~\ref{sec:exp_optimizer}. Importance correction preserves the
conditional expected gradient of the chosen MOPD objective, but AdamW's
second-moment state still depends on the proposal through squared
importance-corrected gradients. Appendix~\ref{app:order_adamw} gives the
state recursion, and the experiments separately measure this order effect.
